%% GitHub Repo Analyzer — Technical Report
%% Compatible with Overleaf (pdfLaTeX)

\documentclass[11pt,a4paper]{article}

\usepackage[margin=2.5cm]{geometry}
\usepackage{booktabs}
\usepackage{longtable}
\usepackage{array}
\usepackage{xcolor}
\usepackage{hyperref}
\usepackage{listings}
\usepackage{enumitem}
\usepackage{titlesec}
\usepackage{parskip}
\usepackage{microtype}
\usepackage{fancyhdr}
\usepackage{graphicx}

% Colour palette
\definecolor{ghblue}{HTML}{0969DA}
\definecolor{codebg}{HTML}{F6F8FA}
\definecolor{tabhead}{HTML}{E8EFF8}
\definecolor{green}{HTML}{1A7F37}
\definecolor{grey}{HTML}{57606A}

% Hyperref setup
\hypersetup{
    colorlinks=true,
    linkcolor=ghblue,
    urlcolor=ghblue,
    pdftitle={GitHub Repo Analyzer — Technical Report},
    pdfauthor={Zahed Riyaz}
}

% Code listing style
\lstset{
    backgroundcolor=\color{codebg},
    basicstyle=\ttfamily\small,
    breaklines=true,
    frame=single,
    framerule=0.4pt,
    rulecolor=\color{grey!40},
    xleftmargin=8pt,
    xrightmargin=8pt,
    aboveskip=6pt,
    belowskip=6pt,
    showstringspaces=false,
    keywordstyle=\color{ghblue},
    commentstyle=\color{grey},
    stringstyle=\color{green}
}

% Section styling
\titleformat{\section}{\large\bfseries\color{ghblue}}{}{0em}{}[\titlerule]
\titleformat{\subsection}{\normalsize\bfseries}{}{0em}{}

% Header / footer
\pagestyle{fancy}
\fancyhf{}
\lhead{\small GitHub Repo Analyzer — Technical Report}
\rhead{\small Zahed Riyaz}
\cfoot{\thepage}

% Column type for wrapped, left-aligned text
\newcolumntype{L}[1]{>{\raggedright\arraybackslash}p{#1}}

% ─────────────────────────────────────────────────────────
\begin{document}

\begin{center}
    {\LARGE\bfseries GitHub Repo Analyzer}\\[4pt]
    {\large Chrome Extension — Technical Report}\\[8pt]
    {\color{grey}\small Manifest V3 $\cdot$ GitHub REST API $\cdot$ Gemini / GPT-4o / Claude / Groq / Ollama $\cdot$ RAG Chat}\\[4pt]
    {\small Zahed Riyaz \quad \textcolor{grey}{|} \quad \today}
\end{center}

\vspace{1em}
\hrule
\vspace{1.5em}

% ─────────────────────────────────────────────────────────
\section{Project Overview}

GitHub Repo Analyzer is a Chrome side-panel extension that surfaces actionable information about any GitHub repository — issues, tech stack, maintainers, contributor health, and an AI chat interface — without a build step. It is built with plain JavaScript, HTML, and CSS under Manifest V3.

\medskip
\textbf{Key capabilities:}
\begin{itemize}[nosep]
    \item Parallel GitHub REST API fetches for repo metadata, issues, languages, contributors, and open PRs.
    \item Tech-stack detection across 50+ tools by scanning dependency files and root directory listings.
    \item Contributor health score (0--100) computed from weighted activity signals.
    \item Retrieval-Augmented Generation (RAG) chat grounded in actual repository source files.
    \item Streaming responses from five AI providers: Gemini, Groq, OpenAI, Anthropic, Ollama.
    \item Per-repo chat history persisted to \texttt{chrome.storage.local}.
\end{itemize}

% ─────────────────────────────────────────────────────────
\section{File Overview}

\begin{longtable}{L{3.4cm} L{1.5cm} L{8.6cm}}
    \toprule
    \rowcolor{tabhead}
    \textbf{File} & \textbf{Type} & \textbf{Purpose} \\
    \midrule
    \endhead

    \bottomrule
    \endfoot

    \texttt{manifest.json} & Config &
        Declares extension metadata, all permissions (\texttt{sidePanel}, \texttt{tabs}, \texttt{storage}) and host permissions for GitHub + five AI provider APIs. \\
    \addlinespace

    \texttt{background.js} & Service Worker &
        Enables the side panel to open on toolbar-icon click via \texttt{sidePanel.setPanelBehavior}; logs extension installation. \\
    \addlinespace

    \texttt{sidepanel.html} & UI Markup &
        Defines the six-tab side-panel layout (Issues, Stack, Maintainers, Contribute, Chat, Settings) and all static DOM scaffolding. \\
    \addlinespace

    \texttt{sidepanel.js} & Main Logic &
        Core 1,800-line script: GitHub API fetching, tab rendering, health scoring, tool detection, RAG context assembly, multi-provider AI streaming, chat history, and settings management. \\
    \addlinespace

    \texttt{options.html} & Settings UI &
        Full-page settings interface with provider selection cards and API-key input fields (mirrors the in-panel settings tab). \\
    \addlinespace

    \texttt{options.js} & Settings Page &
        Handles provider-card selection, API-key persistence to \texttt{chrome.storage.local}, and masked key display on the options page. \\
    \addlinespace

    \texttt{styles.css} & Styling &
        All CSS for the side panel: dark GitHub-like theme, tab layout, chat bubbles, streaming cursor animation, health card, and tool-category grid. \\


\end{longtable}

% ─────────────────────────────────────────────────────────
\section{Key Functions}

\subsection{Initialisation \& Navigation}

\begin{longtable}{L{5.2cm} L{8.3cm}}
    \toprule
    \rowcolor{tabhead}
    \textbf{Function} & \textbf{Description} \\
    \midrule
    \endhead
    \bottomrule
    \endfoot

    \texttt{DOMContentLoaded} listener &
        Extension entry point; loads persisted settings, attaches all UI event listeners, and triggers the first repo load. \\
    \addlinespace

    \texttt{handleRepoRefresh(url)} &
        Detects URL changes; extracts \texttt{owner/repo}; short-circuits on non-repo GitHub pages (e.g.\ trending, gist). \\
    \addlinespace

    \texttt{updateRepoInfo()} &
        Fires five parallel GitHub fetches (repo data, issues, languages, contributors, health + PRs) and stores results in the session cache. \\
\end{longtable}

\subsection{GitHub API Helpers}

\begin{longtable}{L{5.2cm} L{8.3cm}}
    \toprule
    \rowcolor{tabhead}
    \textbf{Function} & \textbf{Description} \\
    \midrule
    \endhead
    \bottomrule
    \endfoot

    \texttt{fetchGitHub(endpoint)} &
        Central API wrapper; adds Bearer token; tracks \texttt{X-RateLimit-Remaining}; updates the rate-limit badge. \\
    \addlinespace

    \texttt{fetchRepoData()} &
        Fetches repo metadata: stars, forks, description, license, fork status. \\
    \addlinespace

    \texttt{fetchIssues()} &
        Fetches open, unassigned issues and filters out pull requests server-side. \\
    \addlinespace

    \texttt{fetchTechStack()} &
        Fetches language byte counts and calls \texttt{detectTools()} to scan dependency files. \\
    \addlinespace

    \texttt{fetchMaintainers()} &
        Fetches top-10 contributors by commit count. \\
    \addlinespace

    \texttt{fetchRepoHealth()} &
        Uses \texttt{Promise.allSettled} to check CONTRIBUTING.md, issue templates, open PRs, and average PR merge time. \\
\end{longtable}

\subsection{Rendering}

\begin{longtable}{L{5.2cm} L{8.3cm}}
    \toprule
    \rowcolor{tabhead}
    \textbf{Function} & \textbf{Description} \\
    \midrule
    \endhead
    \bottomrule
    \endfoot

    \texttt{renderIssues(activeLabel)} &
        Client-side label filter; renders issue cards with coloured label chips. \\
    \addlinespace

    \texttt{renderTechStack(langs, tools)} &
        Displays language percentage bars and groups detected tools by category. \\
    \addlinespace

    \texttt{renderHealthCard(h)} &
        Renders the 0--100 contributor-friendliness score with colour-coded badge and breakdown list. \\
    \addlinespace

    \texttt{renderOpenPRs(prs)} &
        Lists recent open pull requests with author avatar and age. \\
    \addlinespace

    \texttt{renderMarkdown(text)} &
        Lightweight markdown-to-HTML converter: headings, bold, italic, fenced code blocks, ordered/unordered lists, and links. \\
\end{longtable}

\subsection{RAG Chat}

\begin{longtable}{L{5.2cm} L{8.3cm}}
    \toprule
    \rowcolor{tabhead}
    \textbf{Function} & \textbf{Description} \\
    \midrule
    \endhead
    \bottomrule
    \endfoot

    \texttt{handleChat()} &
        Orchestrates a chat turn: calls \texttt{getDeepRepoContext()}, prepends context to the system prompt, builds the last-7-turn history, and streams the AI response. \\
    \addlinespace

    \texttt{getDeepRepoContext()} &
        Fetches up to seven repo files in parallel (README 3,000 chars; others 1,500 chars each), labels each section, joins them, and caches the result for the session. \\
    \addlinespace

    \texttt{appendChatMessage(...)} &
        Renders a chat bubble, adds a timestamp, and wires the Regenerate button. \\
    \addlinespace

    \texttt{saveChatHistory()} &
        Persists the last 50 messages to \texttt{chrome.storage.local} under the key \texttt{chat\_\{owner\}\_\{repo\}}. \\
    \addlinespace

    \texttt{regenerateResponse(...)} &
        Trims history to the point before the last bot turn and re-sends the original query. \\
    \addlinespace

    \texttt{renderChatStarters()} &
        Displays four suggested prompts when the chat history is empty. \\
\end{longtable}

\subsection{AI Provider Routing \& Streaming}

\begin{longtable}{L{5.2cm} L{8.3cm}}
    \toprule
    \rowcolor{tabhead}
    \textbf{Function} & \textbf{Description} \\
    \midrule
    \endhead
    \bottomrule
    \endfoot

    \texttt{callAIStreaming(contents, onChunk)} &
        Top-level router; dispatches to the correct provider streaming function and calls \texttt{onChunk(cumulativeText)} per token. \\
    \addlinespace

    \texttt{callGeminiStreaming(...)} &
        Sends SSE request to Gemini; parses \texttt{candidates[0].content.parts[0].text} events. \\
    \addlinespace

    \texttt{callGroqStreaming(...)} &
        OpenAI-compatible SSE stream from Groq; uses \texttt{geminiToOpenAI()} for format conversion. \\
    \addlinespace

    \texttt{callOpenAIStreaming(...)} &
        Streams from \texttt{gpt-4o-mini}; parses \texttt{choices[0].delta.content} SSE events. \\
    \addlinespace

    \texttt{callAnthropicStreaming(...)} &
        Streams from Claude; parses \texttt{content\_block\_delta / text\_delta} SSE event type. \\
    \addlinespace

    \texttt{callOllamaStreaming(...)} &
        Streams from a local Ollama instance; parses newline-delimited JSON (\texttt{message.content}). \\
    \addlinespace

    \texttt{readOpenAISSEStream(...)} &
        Shared SSE parser for OpenAI/Groq: reads the byte stream, splits on \texttt{data:}, stops at \texttt{[DONE]}. \\
    \addlinespace

    \texttt{geminiToOpenAI(contents)} &
        Converts Gemini message format \texttt{\{role, parts[\{text\}]\}} to OpenAI format \texttt{\{role, content\}}; renames \texttt{"model"} to \texttt{"assistant"}. \\
    \addlinespace

    \texttt{pullOllamaModel(model)} &
        Auto-pulls a missing Ollama model; streams download progress to the typing indicator. \\
\end{longtable}

\subsection{Settings, Scoring \& Utilities}

\begin{longtable}{L{5.2cm} L{8.3cm}}
    \toprule
    \rowcolor{tabhead}
    \textbf{Function} & \textbf{Description} \\
    \midrule
    \endhead
    \bottomrule
    \endfoot

    \texttt{calculateHealthScore(h, repoData)} &
        Returns a 0--100 score from weighted signals: activity (30), CONTRIBUTING.md (20), issue templates (15), PR speed (25), description (5), open issues (5). \\
    \addlinespace

    \texttt{detectTools()} &
        Scans \texttt{package.json}, \texttt{requirements.txt}, and the root directory listing for 50+ tool keywords; groups results by category. \\
    \addlinespace

    \texttt{generateQuickstart()} &
        Sends an AI prompt to write a contributor quickstart guide; streams and caches the result per repo. \\
    \addlinespace

    \texttt{initSettingsTab()} &
        Wires provider-pill clicks, API-key input, Ollama model field, and GitHub token field; persists values to storage. \\
    \addlinespace

    \texttt{maskApiKey(key)} &
        Returns first 6 + \texttt{****} + last 2 characters for safe display. \\
    \addlinespace

    \texttt{escapeHtml(str)} &
        Escapes \texttt{\&}, \texttt{<}, \texttt{>}, \texttt{"} to prevent DOM injection. \\
    \addlinespace

    \texttt{formatNumber(n)} &
        Formats large integers as human-readable strings (e.g.\ \texttt{12\,400} $\rightarrow$ \texttt{12.4k}). \\
    \addlinespace

    \texttt{daysAgo(dateStr)} &
        Converts an ISO date string to a relative label: \texttt{"today"}, \texttt{"3 days ago"}, \texttt{"2mo ago"}, etc. \\
\end{longtable}

% ─────────────────────────────────────────────────────────
\section{RAG Implementation}

\subsection{What is RAG?}

Retrieval-Augmented Generation (RAG) is a technique where relevant external documents are \emph{retrieved} at inference time and \emph{injected} into the LLM prompt as additional context, so the model's answer is grounded in real data rather than relying solely on parametric knowledge from pre-training.

\subsection{Retrieval Layer}

When a user sends a chat message, \texttt{getDeepRepoContext()} fetches the following files from the GitHub Contents API in parallel using \texttt{Promise.allSettled} (missing files are silently skipped):

\medskip
\begin{longtable}{L{4cm} L{2cm} L{7.5cm}}
    \toprule
    \rowcolor{tabhead}
    \textbf{Document} & \textbf{Char Limit} & \textbf{Why It Matters} \\
    \midrule
    \endhead
    \bottomrule
    \endfoot

    \texttt{README} (any ext.) & 3,000 & Project overview, setup, usage — the primary knowledge base. \\
    \addlinespace
    \texttt{CONTRIBUTING.md} & 1,500 & Contribution workflow, coding standards, PR process. \\
    \addlinespace
    \texttt{package.json} & 1,500 & Exact npm dependency names and project scripts. \\
    \addlinespace
    \texttt{requirements.txt} & 1,500 & Python package dependencies. \\
    \addlinespace
    \texttt{pyproject.toml} & 1,500 & Modern Python project config (PEP 517/518). \\
    \addlinespace
    \texttt{Cargo.toml} & 1,500 & Rust crate dependencies. \\
    \addlinespace
    \texttt{Makefile} & 1,500 & Build targets and developer scripts. \\
    \addlinespace
    File tree (\texttt{git/trees/HEAD?recursive=1}) & 5,000 & Paths up to four levels deep, excluding build/vendor directories; reveals project structure. \\
\end{longtable}

Content is returned as base64-encoded strings from the GitHub API and decoded as UTF-8 (\texttt{atob()} to bytes, then \texttt{TextDecoder}), so non-ASCII text survives intact. The entire result is \textbf{cached in the session-scoped \texttt{repoCache}} object so subsequent messages in the same session reuse the context without additional API calls.

\subsection{Augmentation Layer}

The retrieved documents are assembled into a structured context string inside \texttt{getDeepRepoContext()} and then injected into the system prompt inside \texttt{handleChat()}:

\begin{lstlisting}[language=JavaScript]
// Context assembly
parts.push(`=== README ===\n${readmeText.slice(0, 3000)}`);
parts.push(`=== CONTRIBUTING.md ===\n${contribText.slice(0, 1500)}`);
// ... other files ...
parts.push(`=== File Tree ===\n${fileTree}`);
const context = parts.join("\n\n");

// Injection into system prompt
const systemText =
  `You are an expert on the GitHub repository "${owner}/${repo}". ` +
  `Answer questions based on this context:\n\n${context}\n\nUser question: `;

history.push({ role: "user", parts: [{ text: systemText + query }] });
\end{lstlisting}

The last six prior chat turns are prepended before the new message so the model has conversational context as well as document context.

\subsection{Generation Layer}

The augmented prompt is sent to the chosen AI provider via \texttt{callAIStreaming()}, which routes to the correct streaming implementation:

\medskip
\begin{longtable}{L{2.8cm} L{4.4cm} L{2.8cm} L{3.5cm}}
    \toprule
    \rowcolor{tabhead}
    \textbf{Provider} & \textbf{Default Model} & \textbf{Protocol} & \textbf{Context Window} \\
    \midrule
    \endhead
    \bottomrule
    \endfoot

    Gemini  & \texttt{gemini-2.5-flash}           & SSE (\texttt{?alt=sse}) & 1M tokens \\
    Groq    & \texttt{llama-3.3-70b-versatile}    & OpenAI SSE              & 128k tokens \\
    OpenAI  & \texttt{gpt-4o-mini}                & OpenAI SSE              & 128k tokens \\
    Anthropic & \texttt{claude-haiku-4-5-20251001} & Anthropic SSE          & 200k tokens \\
    Ollama  & Configurable (\texttt{llama3.2})    & NDJSON stream           & Model-dependent \\
\end{longtable}

\medskip
Tokens arrive as a byte stream; each provider uses a distinct wire format. All parsers decode the stream incrementally and call \texttt{onChunk(cumulativeText)}, which updates the chat bubble DOM in real time via \texttt{renderMarkdown()}.

\subsection{Why This Works Without a Vector Database}

Traditional RAG pipelines embed documents, store them in a vector database, and retrieve the most semantically similar chunks at query time. This extension takes a simpler approach — \emph{full-document context stuffing} — which is viable here because:

\begin{enumerate}[nosep]
    \item The total size of README + config files typically fits within 5,000--15,000 tokens.
    \item Modern frontier models have 128k--1M token context windows, so the full context always fits.
    \item GitHub repositories have a natural, curated set of authoritative documents (README, CONTRIBUTING, manifest files) that cover the most common contributor questions.
    \item No server infrastructure is required — all retrieval happens client-side via browser \texttt{fetch()}.
\end{enumerate}

The trade-off is that very large READMEs are truncated at 3,000 characters and deep source files are never read. For the contributor-assistance use case this is acceptable.

\subsection{Data Flow Summary}

\begin{lstlisting}
GitHub repo URL (github.com/owner/repo)
        |
        v
  sidepanel.js: chrome.tabs onUpdated / onActivated
  (active tab of the panel's window only)
        |
        v
  sidepanel.js: handleRepoRefresh()
        |
        v
  Parallel GitHub API fetches (metadata, issues, languages,
  contributors, health)  -->  repoCache[owner/repo]
        |
  User sends chat message
        |
        v
  getDeepRepoContext()  -->  7 files fetched in parallel
        |                    (base64-decoded, truncated, labelled)
        v
  Context prepended to system prompt + last 6 turns of history
        |
        v
  callAIStreaming()  -->  routes to Gemini / Groq / OpenAI /
                          Anthropic / Ollama
        |
        v
  SSE / NDJSON stream parsed token-by-token
        |
        v
  onChunk()  -->  renderMarkdown()  -->  DOM update
        |
        v
  saveChatHistory()  -->  chrome.storage.local
\end{lstlisting}

% ─────────────────────────────────────────────────────────
\section{Storage Architecture}

\begin{longtable}{L{3.5cm} L{2.5cm} L{2.5cm} L{5cm}}
    \toprule
    \rowcolor{tabhead}
    \textbf{Store} & \textbf{Scope} & \textbf{Lifetime} & \textbf{Contents} \\
    \midrule
    \endhead
    \bottomrule
    \endfoot

    \texttt{chrome.storage.local} & Persistent & Until cleared & AI provider, API keys, GitHub token, chat history (last 50 messages per repo, keyed \texttt{chat\_owner\_repo}) \\
    \addlinespace
    \texttt{repoCache} (JS object) & Session & Panel reload & Repo metadata, issues, languages, contributors, health, open PRs, quickstart text, RAG context string \\
    \addlinespace
    DOM & Immediate & Render cycle & Chat bubbles, tab content, streaming cursor, rate-limit badge \\
\end{longtable}

% ─────────────────────────────────────────────────────────
\end{document}
